\chapter{Statistical Characterisation of Scar}

\section{AHA-17 Segmentation, Territory Assignment, and Normality Diagnostics}

Scar burden per segment was defined as the fraction of polar-map bins classified as
scar-positive. Territory burden was defined analogously over LAD (segments 1, 2, 7, 8, 13,
14, 17; 4\,448 bins, 54.3\%), RCA (segments 3, 4, 9, 10, 15; 1\,888 bins, 23.0\%), and LCX
(segments 5, 6, 11, 12, 16; 1\,856 bins, 22.7\%). A territory was considered involved when
at least 2\% of its polar bins were scar-positive. This threshold was selected via sensitivity
analysis across candidate values (1\%, 2\%, 3\%, 5\%, 10\%); the full sensitivity curve is
shown in Figure~\ref{fig:app_territory_threshold} in the Appendix. At 1\% the multi-territory
fraction exceeded 60\%, implausibly high for a focal-infarct cohort; thresholds above 3\%
progressively assigned no-scar classification to patients with minor genuine involvement.
The 2\% value represents the earliest stable point at which single-territory classification
reaches its plateau and the no-scar fraction remains low (${\sim}10\%$).

Prior to selecting inferential methods, the distributional properties of all analysis variables
were assessed using quantile--quantile plots and chi-squared goodness-of-fit plots. Of all
variables examined, only age satisfied the normality criterion; all scar-related variables
showed strong zero-inflation and right skewness. All subsequent analyses therefore relied on
non-parametric methods and resampling-based confidence intervals.

Table~\ref{tab:aha17summary} reports prevalence and conditional extent for all 17 AHA segments.

\begin{table}[htbp]
    \centering
    \caption{AHA 17-segment scar burden summary ($n=105$). Segments ordered by ring (basal, mid, apical, apex). Extent = fraction of polar-map bins classified as scar-positive.}
    \label{tab:aha17summary}
    \small
    \begin{tabular}{rllrrrr}
        \toprule
        Seg & Label & Ring & Prev.\ (\%) & Mean (all) & Mean (pos.) & Median (pos.) \\
        \midrule
        1  & Bas ant      & Basal & 30.5 & 0.023 & 0.077 & 0.044 \\
        2  & Bas ant-sep  & Basal & 37.1 & 0.023 & 0.061 & 0.039 \\
        3  & Bas inf-sep  & Basal & 23.8 & 0.014 & 0.059 & 0.030 \\
        4  & Bas inf      & Basal & 33.3 & 0.014 & 0.041 & 0.026 \\
        5  & Bas inf-lat  & Basal & 44.8 & 0.032 & 0.071 & 0.033 \\
        6  & Bas ant-lat  & Basal & 39.0 & 0.032 & 0.083 & 0.063 \\
        \midrule
        7  & Mid ant      & Mid   & 34.3 & 0.042 & 0.121 & 0.048 \\
        8  & Mid ant-sep  & Mid   & 48.6 & 0.057 & 0.117 & 0.063 \\
        9  & Mid inf-sep  & Mid   & 52.4 & 0.060 & 0.114 & 0.060 \\
        10 & Mid inf      & Mid   & 61.0 & 0.088 & 0.145 & 0.080 \\
        11 & Mid inf-lat  & Mid   & 56.2 & 0.100 & 0.178 & 0.110 \\
        12 & Mid ant-lat  & Mid   & 46.7 & 0.063 & 0.136 & 0.071 \\
        \midrule
        13 & Api ant      & Apical & 41.0 & 0.060 & 0.148 & 0.066 \\
        14 & Api sep      & Apical & 41.0 & 0.078 & 0.190 & 0.098 \\
        15 & Api inf      & Apical & 56.2 & 0.058 & 0.103 & 0.059 \\
        16 & Api lat      & Apical & 49.5 & 0.051 & 0.103 & 0.059 \\
        \midrule
        17 & Apex         & Apex  & 53.3 & 0.052 & 0.097 & 0.024 \\
        \bottomrule
    \end{tabular}
\end{table}

\section{Within-Patient Spatial Tests and Bootstrap Confidence Intervals}

Whether scar burden varies systematically across AHA segments within the same patient
was tested using the Friedman test, a non-parametric analogue of one-way
repeated-measures ANOVA. Each of the 105 patients contributed a vector of 17 paired
segment-burden values. Pairwise post-hoc comparisons were conducted using Wilcoxon
signed-rank tests with Bonferroni correction for the 136 pairwise comparisons.

Population-level confidence intervals for the mean scar burden in each AHA segment were
estimated via patient-level bootstrap resampling. In each of 2\,000 iterations, 105 patients
were sampled with replacement---preserving each patient's full 17-segment vector---and the
mean burden per segment was recomputed. The 2.5th and 97.5th percentiles defined the
95\% confidence interval. The primary purpose of this procedure is not to establish
statistical significance of individual segment means, but to quantify the uncertainty around
population estimates for the atlas and archetype maps. Because the underlying distributions
are zero-inflated and heavily skewed, symmetric normal-theory intervals would systematically
misrepresent the sampling distribution: they produce symmetric error bars around means
that are, in reality, bounded below by zero and pulled by long positive tails. The bootstrap
intervals are asymmetric by construction and therefore faithfully reflect the shape of the
sampling distribution, providing honest uncertainty bounds for the population scar atlas.

\section{Age and Sex Associations}

Association between age and segment-level scar burden was assessed using Spearman
rank correlation in the 48 patients with available age labels (mean age $63.0 \pm 11.5$
years). Bonferroni correction was applied across the 17 segments, yielding an adjusted
significance threshold of $\alpha = 0.05/17 = 0.0029$. Sex-stratified analysis was limited by
severe label imbalance (38 male, 10 female among patients with sex information).

\section{Dimensionality Reduction and Correlation Structure}

PCA was applied to the 17-dimensional scar-burden vectors after $z$-score standardisation
(zero mean, unit variance per segment). The aim was to identify the dominant spatial modes
of scar variability and assess the intrinsic dimensionality of the scar-pattern space.
Component loadings were mapped back onto the AHA bullseye representation to visualise
the spatial scar phenotype associated with each PC.

Two correlation analyses complement the PCA. First, the inter-segment correlation matrix
computes Spearman rank correlations between all $17 \times 17$ pairs of AHA segment
burdens across patients. Second, the territory burden correlation matrix
computes the $3 \times 3$ Spearman correlation matrix between LAD, RCA, and LCX burden.
This determines whether territory burdens are effectively independent dimensions---weak
correlations would support treating territories as orthogonal axes of scar phenotype---or
whether multi-territory involvement is the norm, implying that a single ``core zone burden''
scalar captures most of the territory-level information.

\section{Population Variability and Archetype Maps}

Population variability maps were computed across the polar maps. Archetype maps were
constructed by grouping patients according to their dominant coronary territory and
computing the mean polar map within each group. Ring-level and territory-level coverage
summaries, bootstrap confidence intervals, and variability maps were recomputed using the
area-based and coverage projections.

\section{Transmurality: Continuous Vertex-Depth Computation}

The analyses above characterise scar extent on the LV surface but do not capture how
deeply scar penetrates through the ventricular wall. Transmurality---the fraction of wall
thickness affected by scar---is clinically important because transmural scar has different
electrophysiological consequences than subendocardial scar.

An initial approach used the nine layer-specific scar meshes exported by ADAS 3D
(Layer\_10 through Layer\_90), projecting each layer mesh onto the $64 \times 128$ polar
grid and computing onset and elongation as the shallowest and span of scar-containing
layers. However, the discrete layer resolution proved too coarse: the resulting transmurality
values were nearly uniform across the cohort, with insufficient variation to differentiate scar
profiles meaningfully.

To overcome this limitation, a continuous vertex-depth approach was developed. For each
patient, the endocardial and epicardial surfaces of the LV wall are loaded, and a cKDTree
spatial index is built for each. Every vertex of the core-zone scar mesh is then located in 3D
space, and its Euclidean distances to the nearest point on the endocardial surface ($d_\text{endo}$)
and the epicardial surface ($d_\text{epi}$) are computed. The fractional wall depth of that vertex
is defined as:
\begin{equation}
    \text{depth\_frac} = \frac{d_\text{endo}}{d_\text{endo} + d_\text{epi}}.
    \label{eq:depthfrac}
\end{equation}
This yields a continuous value in $[0, 1]$, where 0 represents the endocardial wall and 1 the
epicardial wall. Each scar vertex is then projected onto the same $64 \times 128$ polar grid
used for the binary scar maps. Within each polar bin, the minimum depth fraction across all
scar vertices falling into that bin defines the \emph{onset} (the shallowest wall depth at which
scar begins), and the maximum depth fraction defines the \emph{outermost depth}.
Transmurality per bin is then computed as:
\begin{equation}
    \text{transmurality} = \text{outermost\_depth} - \text{onset}.
    \label{eq:transmurality}
\end{equation}

These bin-level values are aggregated at three levels (segment, territory, and fragment) over
scar-positive bins only; the aggregation procedure is detailed in
Appendix~\ref{app:transmuralagg}. This decomposition distinguishes three clinically
relevant scar profiles: thin subendocardial (low onset, low transmurality), transmural (low
onset, high transmurality), and midwall (high onset, variable transmurality).

\section{Fragment-Level Characterisation}
\label{sec:fragments}

Connected-component labelling was applied to each patient's 3D core-zone mesh to identify
spatially distinct scar fragments (minimum 50 vertices). Each fragment is projected to the
polar grid and assigned to a ring zone (apex, apical, mid, basal) by its centroid row and to a
coronary territory by plurality vote across its polar bins. For each fragment, the mean onset,
mean outermost depth, and mean transmurality are computed from the continuous
vertex-depth grids over the fragment's polar footprint. Fragment size (number of polar bins),
ring position, and territory assignment complete the per-fragment descriptor set.

\section{Within-Patient Consistency Between Extent and Transmurality}

For each patient, the 17-element segment-level extent vector (fraction of bins that are
scar-positive per AHA segment) is correlated with the 17-element per-segment mean
transmurality vector (mean transmurality over scar-positive bins in each segment) using
Spearman rank correlation. Values near zero indicate that extent and transmurality carry
independent information about the scar phenotype---a patient may have extensive but
shallow scar, or focal but deeply penetrating scar---and that both descriptors are needed.
Values near $+1$ indicate redundancy, where one measure largely substitutes for the other.

\section{Transmurality Gradient Within Fragments}

For each 2D polar-map fragment with at least 20 bins, each bin is classified as \emph{core}
(three or more scar-positive 4-connected neighbours) or \emph{edge} (fewer than three).
The fragment gradient is then defined as the difference between the mean transmurality of
core bins and the mean transmurality of edge bins. A positive population-mean gradient
indicates that scar has a dense-core / thin-edge transmural structure, which is relevant for
arrhythmogenesis because the gradient zone---where transmurality transitions from thick to
thin---is precisely where conduction slows and re-entry circuits may form.

\section{Geometric Shape Descriptors}

Four additional per-patient descriptors capture aspects of scar geometry not encoded in
segment-level or transmurality analyses. \textbf{Circumferential extent} is the fraction of
angular columns containing at least one scar-positive bin per ring level (apex, apical, mid,
basal), quantifying tangential spread around the ventricular wall. \textbf{Scar compactness}
is $4\pi \cdot \text{area}/\text{perimeter}^2$ of the largest 2D connected component (1 for
a perfect circle, lower for elongated or irregular shapes), supplemented by aspect ratio and
fragment count; these metrics are independent of total scar burden and capture whether scar
forms a single compact lesion or a fragmented pattern. \textbf{Endocardial--epicardial
dominance} classifies each scar-positive bin by onset depth: endocardial ($\leq 30\%$),
midwall (30--60\%), or epicardial ($> 60\%$ of wall thickness), distinguishing ischaemic
from non-ischaemic onset patterns. \textbf{Apex-to-base gradient} fits a first-order
polynomial to row-level mean burden ($s \in [0,1]$, apex to base), with Spearman
correlation as a non-parametric summary of longitudinal scar organisation.
